\documentclass[main.tex]{subfiles}
\begin{document}

\section{Exercise 6 – Input capacitance of the cascode amplifier (Figure~1.3)}

\subsection*{Small-signal model}
The cascode adds transistor \(M_3\) in a common-gate configuration, holding the node \(v_X\)
(drains of \(M_1\)) at nearly constant AC voltage, thereby reducing the Miller effect.
The small-signal model is shown in Figure~\ref{fig:ss_ex6}.

\begin{figure}[H]
    \centering
    \includegraphics[width=0.5\linewidth]{images/small_signal_model_ex6.png}
    \caption{Small-signal circuit of the cascode amplifier.}
    \label{fig:ss_ex6}
\end{figure}

\subsection*{KCL and calculation}
Applying KCL at node \(v_X\):
\[
0=g_{m3}v_x+g_{m1}v_{in}
\quad\Rightarrow\quad
A_X=\frac{v_x}{v_{in}}=-\frac{g_{m1}}{g_{m3}}.
\]
For matched transistors (\(g_{m1}\!\approx\!g_{m3}\)), \(A_X\!\approx\!-1\), giving
\[
C_{\text{in}}=C_1+C_2(1-A_X)=C_1+C_2(1+1)=0.1+0.1=0.2~\text{pF}.
\]

\(\boxed{C_{\text{in}}=0.2~\text{pF}}\).

\subsection*{Result}
The cascode stage greatly reduces the Miller effect compared to the common-source amplifier,
lowering the effective input capacitance from \(1.31~\text{pF}\) to around \(0.2~\text{pF}\).

\end{document}
